\documentclass[11pt,a4paper]{article}

\usepackage[utf8]{inputenc}
\usepackage{amsmath,amssymb}
\usepackage{booktabs,tabularx}
\usepackage{enumitem}
\usepackage{geometry}
\usepackage{hyperref}
\usepackage{xcolor}
\usepackage{listings}
\usepackage{fancyhdr}
\usepackage{titlesec}

\geometry{margin=2.3cm,top=2.5cm,bottom=2.5cm}

\hypersetup{
  colorlinks=true,
  linkcolor=blue!60!black,
  urlcolor=blue!60!black,
  pdftitle={ISLS v3.2 --- Full Pipeline Wiring},
  pdfauthor={Sebastian Klemm},
}

\lstset{
  basicstyle=\ttfamily\scriptsize,
  breaklines=true,
  frame=single,
  backgroundcolor=\color{gray!8},
  keywordstyle=\color{blue!70!black}\bfseries,
  commentstyle=\color{green!50!black}\itshape,
  numbers=left,
  numberstyle=\tiny\color{gray},
  numbersep=5pt,
  tabsize=4,
  showstringspaces=false,
}

\pagestyle{fancy}
\fancyhf{}
\fancyhead[L]{\small ISLS v3.2 --- Full Pipeline Wiring}
\fancyhead[R]{\small Klemm, 2026}
\fancyfoot[C]{\thepage}

\begin{document}

\begin{titlepage}
\begin{center}
\vspace*{1.5cm}
{\Huge\bfseries ISLS v3.2}\\[0.4cm]
{\LARGE\bfseries Full Pipeline Wiring}\\[0.3cm]
{\LARGE\bfseries Chat $\to$ Norms $\to$ Forge $\to$ App}\\[0.3cm]
{\LARGE\bfseries Multi-Domain Proof \& Live Studio}\\[1.2cm]
{\Large Specification v1.0}\\[1.5cm]
{\large Sebastian Klemm}\\[0.3cm]
{\small March 2026}\\[2cm]

\begin{tabular}{ll}
\textbf{Repository:} & \texttt{graphity/isls} \\
\textbf{Builds on:} & v3.1 (Template-Free LLM Generation) \\
\textbf{Modified crates:} & \texttt{isls-gateway}, \texttt{isls-agent}, \\
 & \texttt{isls-studio}, \texttt{isls-chat}, \\
 & \texttt{isls-forge-llm}, \texttt{isls-cli} \\
\textbf{No new crates.} & Everything exists. This spec \emph{wires} it. \\
\textbf{Goal:} & Type in the browser. Get a working app. \\
\end{tabular}

\vfill
{\small
All the pieces are built. None of them talk to each other.\\
This spec connects: Studio chat $\to$ intent recognition $\to$\\
norm composition $\to$ template-free LLM forge $\to$\\
compiled output $\to$ Studio dashboard with live progress.\\[0.3cm]
Three domains verified: Warehouse, E-Commerce, Project Tracker.\\
One unified experience: type what you want, get what you need.
}
\end{center}
\end{titlepage}

\tableofcontents
\newpage

% ============================================================
\section{What Exists But Is Not Connected}

\begin{center}
\small
\begin{tabular}{lllp{4.5cm}}
\toprule
\textbf{Crate} & \textbf{Status} & \textbf{Has} & \textbf{Missing} \\
\midrule
isls-chat & Built & Intent types, keyword extraction & Not called by gateway \\
isls-norms & Built & 24 norms, composition, learning & Not called by forge pipeline \\
isls-forge-llm & Built & LLM gen, cargo check, retry & Only called via CLI, not gateway \\
isls-gateway & Built & 9 endpoints, serve command & POST /api/chat returns static response \\
isls-studio & Built & HTML, chat panel, norm explorer & Chat sends messages but gets no real response \\
isls-agent & Built & Pipeline executor & Not wired to forge-llm \\
\bottomrule
\end{tabular}
\end{center}

The fix is not building new things. It is connecting existing things.

% ============================================================
\section{The Connected Pipeline}

\subsection{Flow: Chat Message to Running App}

\begin{lstlisting}[title=Complete Flow]
USER types in Studio: "Warehouse management system with products,
                       orders, suppliers, and inventory tracking"
  |
  | POST /api/chat { message: "..." }
  v
GATEWAY receives message
  |
  | calls isls-chat::extract_intent()
  v
CHAT extracts: {
    intent: CreateApplication,
    entities: ["product", "order", "supplier", "warehouse",
               "order_item", "stock_movement", "user"],
    features: ["inventory_tracking", "order_management"],
    roles: ["admin", "operator", "viewer"]
  }
  |
  | calls isls-norms::NormRegistry::match_description()
  v
NORMS activates: [
    NORM-0042 x7 (CRUD for each entity),
    NORM-0088 (JWT Auth),
    NORM-0112 (Inventory Tracking),
    NORM-0120 (Order State Machine),
    NORM-0100 (Pagination),
    NORM-0101 (Error System),
    NORM-0210 (Config),
    NORM-0220 (Health Check),
    NORM-0230 (Docker Deploy)
  ]
  |
  | calls isls-norms::compose_norms()
  v
COMPOSED PLAN: {
    entities: 7 with full field definitions,
    services: 7 x CRUD + adjust_stock + fulfill_order,
    api: 35+ endpoints,
    frontend: 7 pages + dashboard + login
  }
  |
  | GATEWAY returns ChatResponse to Studio
  v
STUDIO shows:
  "I'll build a warehouse management system with:
   - 7 data types (Product, Order, Supplier, ...)
   - Inventory tracking with stock movements
   - Order state machine (pending -> shipped -> delivered)
   - JWT authentication with 3 roles
   - Dashboard with KPIs

   Estimated: ~35 files, ~4,000 LOC
   [Generate]  [Customize]"
  |
  | USER clicks [Generate]
  | POST /api/projects/forge { plan_id: "...", settings: {...} }
  v
GATEWAY spawns generation thread
  |
  | calls isls-forge-llm::LlmForge::generate()
  v
FORGE generates file by file:
  Layer 0: static files (Cargo.toml, Docker, etc.)
  Layer 1: errors.rs, config.rs, pagination.rs
  Layer 2: auth.rs, models/user.rs
  Layer 3: models/*.rs
  ...
  |
  | After EACH file: POST progress to STUDIO via SSE/polling
  v
STUDIO shows live progress:
  [####-----------] 12/35 files
  Layer 3: Generating models/order.rs... (attempt 1)
  Tokens: 12,450 | Compile checks: 10 | Retries: 1
  |
  | When complete:
  v
STUDIO shows:
  "Generation complete!
   35 files, 4,120 LOC, 28,900 tokens ($0.004)
   Compile: PASS (2 retries on services/order.rs)
   Evidence chain: valid (SHA-256)

   [Download ZIP]  [View Files]  [docker-compose up instructions]"
\end{lstlisting}

% ============================================================
\section{Gateway Wiring}

\subsection{POST /api/chat --- Real Implementation}

\begin{lstlisting}[language=C]
async fn handle_chat(
    State(state): State<AppState>,
    Json(body): Json<ChatRequest>,
) -> Result<Json<ChatResponse>, StatusCode> {
    let message = &body.message;

    // 1. Extract intent
    let intent = if let Some(oracle) = &state.oracle {
        isls_chat::extract_intent(oracle.as_ref(), message)?
    } else {
        isls_chat::extract_intent_keywords(message)
    };

    // 2. Based on intent type, build response
    match intent.intent_type {
        IntentType::CreateApplication => {
            // Match norms
            let activated = state.norm_registry.match_from_intent(&intent);
            // Compose
            let plan = isls_norms::compose_norms(&activated, &HashMap::new())?;
            // Estimate
            let est_files = plan.estimated_file_count();
            let est_loc = plan.estimated_loc();
            // Store plan for later [Generate] click
            let plan_id = state.store_plan(plan.clone()).await;

            Ok(Json(ChatResponse {
                message: format_plan_description(&plan, &activated),
                intent: "create_application".into(),
                norms_activated: activated.iter()
                    .map(|n| n.norm.id.clone()).collect(),
                estimated_files: est_files,
                estimated_loc: est_loc,
                plan_id: Some(plan_id),
                action_buttons: vec![
                    ActionButton {
                        label: "Generate".into(),
                        action: format!("/api/projects/forge"),
                        style: "primary".into(),
                    },
                    ActionButton {
                        label: "Customize".into(),
                        action: "/api/chat".into(),
                        style: "secondary".into(),
                    },
                ],
            }))
        }
        IntentType::AddField => {
            // Handle amendment to existing project
            let ops = isls_chat::intent_to_norm_ops(&intent, &state.norm_registry);
            let affected = ops.iter()
                .flat_map(|op| isls_chat::affected_files(op))
                .collect::<Vec<_>>();
            Ok(Json(ChatResponse {
                message: format!("I'll add {} to {}. This affects {} files.",
                    intent.fields.first().map(|f| &f.field_name).unwrap_or(&"?".into()),
                    intent.fields.first().map(|f| &f.entity).unwrap_or(&"?".into()),
                    affected.len()),
                action_buttons: vec![ActionButton {
                    label: "Apply Changes".into(),
                    action: "/api/projects/amend".into(),
                    style: "primary".into(),
                }],
                ..Default::default()
            }))
        }
        IntentType::Help => {
            Ok(Json(ChatResponse {
                message: help_text(),
                ..Default::default()
            }))
        }
        _ => {
            Ok(Json(ChatResponse {
                message: "I'm not sure what you want. Try describing \
                         the app you need, or say 'help'.".into(),
                ..Default::default()
            }))
        }
    }
}

fn format_plan_description(
    plan: &ComposedPlan,
    activated: &[ActivatedNorm],
) -> String {
    let mut msg = String::new();
    msg.push_str(&format!("I'll build a system with:\n"));
    // List entities
    let entity_count = plan.models.len();
    msg.push_str(&format!("- {} data types", entity_count));
    if entity_count <= 10 {
        let names: Vec<_> = plan.models.iter()
            .map(|m| &m.struct_name).collect();
        msg.push_str(&format!(" ({})", names.join(", ")));
    }
    msg.push_str("\n");
    // List key features from norms
    for norm in activated {
        match norm.norm.id.as_str() {
            "ISLS-NORM-0088" => msg.push_str("- JWT authentication with role-based access\n"),
            "ISLS-NORM-0112" => msg.push_str("- Inventory tracking with stock movements\n"),
            "ISLS-NORM-0120" => msg.push_str("- Status workflow with state machine\n"),
            "ISLS-NORM-0130" => msg.push_str("- Full-text search\n"),
            "ISLS-NORM-0180" => msg.push_str("- Dashboard with KPI cards\n"),
            "ISLS-NORM-0190" => msg.push_str("- Charts and data visualization\n"),
            _ => {}
        }
    }
    msg.push_str(&format!("\nEstimated: ~{} files, ~{} lines of code\n",
        plan.estimated_file_count(), plan.estimated_loc()));
    msg
}
\end{lstlisting}

\subsection{POST /api/projects/forge --- Real Generation}

\begin{lstlisting}[language=C]
async fn handle_forge(
    State(state): State<AppState>,
    Json(body): Json<ForgeRequest>,
) -> Result<Json<ForgeResponse>, StatusCode> {
    // 1. Retrieve stored plan
    let plan = state.get_plan(&body.plan_id).await
        .ok_or(StatusCode::NOT_FOUND)?;

    // 2. Build AppSpec from ComposedPlan
    let app_spec = plan_to_app_spec(&plan, &body.settings)?;

    // 3. Build ForgePlan from AppSpec
    let forge_plan = ForgePlan::from_domain(&app_spec);

    // 4. Create output directory
    let output_dir = state.projects_dir.join(&app_spec.name);
    std::fs::create_dir_all(&output_dir)?;

    // 5. Create LlmForge
    let oracle: Box<dyn Oracle> = if let Some(key) = &body.api_key {
        Box::new(OpenAiOracle::new(key.clone()))
    } else {
        Box::new(MockOracle)
    };

    let mut forge = LlmForge::new(
        oracle,
        forge_plan,
        app_spec.clone(),
        output_dir.clone(),
    );

    // 6. Store project metadata
    let project_id = state.create_project(&app_spec, &output_dir).await;

    // 7. Generate (this takes 1-5 minutes with LLM)
    //    In production: spawn a thread + progress polling.
    //    For v3.2: blocking is acceptable.
    let result = forge.generate(&app_spec)?;

    // 8. Store result
    state.update_project_result(&project_id, &result).await;

    // 9. Run norm learning
    state.norm_registry.lock().await
        .observe_and_learn_from_forge(&result);

    Ok(Json(ForgeResponse {
        project_id,
        files_generated: result.len(),
        total_loc: result.iter().map(|f| f.content.lines().count()).sum(),
        total_tokens: result.iter().map(|f| f.tokens_used).sum(),
        compile_status: "pass".into(),
        output_dir: output_dir.to_string_lossy().into(),
    }))
}
\end{lstlisting}

\subsection{GET /api/projects/:id/files --- File Browser}

\begin{lstlisting}[language=C]
async fn list_project_files(
    State(state): State<AppState>,
    Path(project_id): Path<String>,
) -> Result<Json<Vec<FileInfo>>, StatusCode> {
    let project = state.get_project(&project_id).await
        .ok_or(StatusCode::NOT_FOUND)?;
    let files = walk_directory(&project.output_dir)?;
    Ok(Json(files))
}

async fn get_project_file(
    State(state): State<AppState>,
    Path((project_id, file_path)): Path<(String, String)>,
) -> Result<String, StatusCode> {
    let project = state.get_project(&project_id).await
        .ok_or(StatusCode::NOT_FOUND)?;
    let full_path = project.output_dir.join(&file_path);
    std::fs::read_to_string(&full_path)
        .map_err(|_| StatusCode::NOT_FOUND)
}
\end{lstlisting}

\subsection{AppState}

\begin{lstlisting}[language=C]
pub struct AppState {
    pub oracle: Option<Box<dyn Oracle + Send + Sync>>,
    pub norm_registry: Arc<Mutex<NormRegistry>>,
    pub projects: Arc<Mutex<Vec<ProjectRecord>>>,
    pub pending_plans: Arc<Mutex<HashMap<String, ComposedPlan>>>,
    pub projects_dir: PathBuf,      // ~/.isls/projects/
    pub crystal_registry: Arc<Mutex<CrystalRegistry>>,
}

pub struct ProjectRecord {
    pub id: String,
    pub name: String,
    pub description: String,
    pub output_dir: PathBuf,
    pub created_at: String,
    pub status: ProjectStatus,
    pub stats: Option<ForgeStats>,
    pub norms_used: Vec<String>,
}

pub enum ProjectStatus {
    Planning,       // plan created, not yet generated
    Generating,     // LLM forge in progress
    Completed,      // generation done, compilable
    Failed(String), // generation failed
}
\end{lstlisting}

% ============================================================
\section{Studio Dashboard --- Live Wiring}

\subsection{Chat Panel Wiring}

The Studio chat panel already has the HTML. Wire it to real API calls:

\begin{lstlisting}[language=Java,title=Studio JavaScript (in studio.html)]
// Chat submit handler
async function sendMessage() {
    const input = document.getElementById('chat-input');
    const message = input.value.trim();
    if (!message) return;

    // Show user message
    appendMessage('user', message);
    input.value = '';

    // Call API
    try {
        const resp = await fetch('/api/chat', {
            method: 'POST',
            headers: { 'Content-Type': 'application/json' },
            body: JSON.stringify({ message }),
        });
        const data = await resp.json();

        // Show system response
        appendMessage('system', data.message);

        // Show action buttons
        if (data.action_buttons && data.action_buttons.length > 0) {
            const btnContainer = document.createElement('div');
            btnContainer.className = 'action-buttons';
            for (const btn of data.action_buttons) {
                const button = document.createElement('button');
                button.textContent = btn.label;
                button.className = `btn btn-${btn.style}`;
                button.onclick = () => handleAction(btn, data.plan_id);
                btnContainer.appendChild(button);
            }
            document.getElementById('chat-messages').appendChild(btnContainer);
        }

        // Update norm explorer if norms were activated
        if (data.norms_activated && data.norms_activated.length > 0) {
            updateNormExplorer(data.norms_activated);
        }
    } catch (err) {
        appendMessage('system', 'Error: ' + err.message);
    }
}

// Handle action button click
async function handleAction(btn, planId) {
    if (btn.label === 'Generate') {
        appendMessage('system', 'Starting generation...');
        showProgressBar();

        const settings = getForgeSettings(); // API key input, model selector
        try {
            const resp = await fetch('/api/projects/forge', {
                method: 'POST',
                headers: { 'Content-Type': 'application/json' },
                body: JSON.stringify({
                    plan_id: planId,
                    api_key: settings.apiKey || null,
                    settings: settings,
                }),
            });
            const result = await resp.json();

            hideProgressBar();
            appendMessage('system',
                `Generation complete!\n` +
                `${result.files_generated} files, ${result.total_loc} LOC\n` +
                `Tokens: ${result.total_tokens}\n` +
                `Status: ${result.compile_status}`
            );

            // Refresh projects list
            loadProjects();
        } catch (err) {
            hideProgressBar();
            appendMessage('system', 'Generation failed: ' + err.message);
        }
    }
}
\end{lstlisting}

\subsection{Projects Panel Wiring}

\begin{lstlisting}[language=Java,title=Projects Panel]
async function loadProjects() {
    const resp = await fetch('/api/projects');
    const projects = await resp.json();

    const container = document.getElementById('projects-list');
    container.innerHTML = '';

    for (const p of projects) {
        const card = document.createElement('div');
        card.className = 'project-card';
        card.innerHTML = `
            <h3>${p.name}</h3>
            <div class="project-meta">
                <span class="status status-${p.status.toLowerCase()}">${p.status}</span>
                <span>${p.stats ? p.stats.files_generated + ' files' : ''}</span>
                <span>${p.stats ? p.stats.total_loc + ' LOC' : ''}</span>
                <span>${p.created_at}</span>
            </div>
            <div class="project-norms">
                ${p.norms_used.map(n => `<span class="norm-badge">${n}</span>`).join('')}
            </div>
            <div class="project-actions">
                <button onclick="viewFiles('${p.id}')" class="btn btn-secondary">Files</button>
                <button onclick="downloadZip('${p.id}')" class="btn btn-secondary">Download</button>
            </div>
        `;
        container.appendChild(card);
    }
}

async function viewFiles(projectId) {
    const resp = await fetch(`/api/projects/${projectId}/files`);
    const files = await resp.json();
    showFileTree(files, projectId);
}

async function showFileContent(projectId, filePath) {
    const resp = await fetch(`/api/projects/${projectId}/files/${filePath}`);
    const content = await resp.text();
    document.getElementById('file-viewer').textContent = content;
    document.getElementById('file-viewer-name').textContent = filePath;
    document.getElementById('file-viewer-panel').style.display = 'block';
}
\end{lstlisting}

\subsection{Norm Explorer Wiring}

\begin{lstlisting}[language=Java,title=Norm Explorer]
async function loadNorms() {
    const resp = await fetch('/api/norms');
    const norms = await resp.json();

    const container = document.getElementById('norms-list');
    container.innerHTML = '';

    for (const norm of norms) {
        const el = document.createElement('div');
        const levelClass = norm.level.toLowerCase();
        const sourceClass = norm.source === 'auto' ? 'auto' : 'builtin';
        el.className = `norm-item norm-${levelClass} norm-${sourceClass}`;
        el.innerHTML = `
            <span class="norm-id">${norm.id}</span>
            <span class="norm-name">${norm.name}</span>
            <span class="norm-level">${norm.level}</span>
            ${norm.source === 'auto' ?
                `<span class="norm-obs">${norm.observations} obs, ${norm.domains} domains</span>` : ''}
        `;
        container.appendChild(el);
    }
}

function updateNormExplorer(activatedIds) {
    // Highlight activated norms in the explorer
    document.querySelectorAll('.norm-item').forEach(el => {
        el.classList.remove('activated');
    });
    for (const id of activatedIds) {
        const el = document.querySelector(`[data-norm-id="${id}"]`);
        if (el) el.classList.add('activated');
    }
}
\end{lstlisting}

\subsection{Settings Panel}

Add a settings area where the user enters their API key:

\begin{lstlisting}[language=Java,title=Settings]
function getForgeSettings() {
    return {
        apiKey: document.getElementById('api-key-input').value || null,
        model: document.getElementById('model-select').value || 'gpt-4o-mini',
    };
}
\end{lstlisting}

The API key is stored in \texttt{localStorage} (browser-side only).
Never sent to the ISLS server except when calling
\texttt{/api/projects/forge}.

% ============================================================
\section{Multi-Domain: ForgePlan from Norms}

\subsection{The Problem}

v3.1's \texttt{ForgePlan::warehouse\_default()} is hardcoded for
warehouse. To support any domain, the ForgePlan must be built
from the ComposedPlan's entities.

\subsection{ForgePlan::from\_composed\_plan()}

\begin{lstlisting}[language=C]
impl ForgePlan {
    /// Build a ForgePlan from a ComposedPlan (from norm composition).
    /// This replaces warehouse_default() for norm-driven generation.
    pub fn from_composed_plan(plan: &ComposedPlan, app_spec: &AppSpec) -> Self {
        let mut files = Vec::new();

        // Layer 0: Static files
        files.push(FileSpec::static_file("Cargo.toml"));
        files.push(FileSpec::static_file("docker-compose.yml"));
        files.push(FileSpec::static_file("Dockerfile"));
        files.push(FileSpec::static_file(".env.example"));
        files.push(FileSpec::static_file(".gitignore"));

        // Layer 1: Foundation
        files.push(FileSpec::llm_file("src/errors.rs", "Error types", 1));
        files.push(FileSpec::llm_file("src/config.rs", "App configuration", 1));
        files.push(FileSpec::llm_file("src/pagination.rs", "Pagination types", 1));

        // Layer 2: Auth (always, since NORM-0088 is always activated)
        // Find the User-like entity (User, Customer, Account, etc.)
        let user_entity = find_user_entity(&plan.models)
            .unwrap_or_else(|| standard_user_entity());
        files.push(FileSpec::llm_file(
            &format!("src/models/{}.rs", user_entity.snake_name()), "User model", 2));
        files.push(FileSpec::llm_file("src/auth.rs", "JWT authentication", 2));

        // Layer 3: All other models (from norm composition)
        for model in &plan.models {
            if model.struct_name == user_entity.struct_name { continue; } // already in layer 2
            files.push(FileSpec::llm_file(
                &format!("src/models/{}.rs", model.snake_name()),
                &format!("{} model", model.struct_name), 3));
        }
        files.push(FileSpec::llm_file("src/models/mod.rs", "Model declarations", 3));

        // Layer 4: Database
        files.push(FileSpec::llm_file(
            "database/migrations/001_initial.sql", "Database schema", 4));
        for model in &plan.models {
            files.push(FileSpec::llm_file(
                &format!("src/database/{}_queries.rs", model.snake_name()),
                &format!("{} queries", model.struct_name), 4));
        }
        files.push(FileSpec::llm_file("src/database/pool.rs", "DB pool", 4));
        files.push(FileSpec::llm_file("src/database/mod.rs", "DB declarations", 4));

        // Layer 5: Services
        for model in &plan.models {
            files.push(FileSpec::llm_file(
                &format!("src/services/{}.rs", model.snake_name()),
                &format!("{} business logic", model.struct_name), 5));
        }
        files.push(FileSpec::llm_file("src/services/mod.rs", "Service declarations", 5));

        // Layer 6: API
        for model in &plan.models {
            files.push(FileSpec::llm_file(
                &format!("src/api/{}.rs", model.snake_name()),
                &format!("{} endpoints", model.struct_name), 6));
        }
        files.push(FileSpec::llm_file("src/api/auth_routes.rs", "Auth endpoints", 6));
        files.push(FileSpec::llm_file("src/api/mod.rs", "API declarations", 6));

        // Layer 7: Main
        files.push(FileSpec::llm_file("src/main.rs", "Application entry", 7));

        // Layer 8: Frontend
        files.push(FileSpec::llm_file("frontend/index.html", "SPA shell", 8));
        files.push(FileSpec::llm_file("frontend/style.css", "Styles", 8));
        files.push(FileSpec::llm_file("frontend/app.js", "App logic", 8));

        // Layer 9: Tests
        files.push(FileSpec::llm_file("tests/api_tests.rs", "API tests", 9));

        ForgePlan {
            app_name: app_spec.name.clone(),
            files,
            entities: plan.models.iter().map(|m| model_to_entity(m)).collect(),
        }
    }

    /// Find a User-like entity in the models.
    fn find_user_entity(models: &[ModelArtifact]) -> Option<EntityDef> {
        // Look for: User, Customer, Account, Member, Employee
        let user_names = ["User", "Customer", "Account", "Member", "Employee"];
        for name in &user_names {
            if let Some(model) = models.iter().find(|m| m.struct_name == *name) {
                return Some(model_to_entity(model));
            }
        }
        None
    }
}
\end{lstlisting}

\subsection{Standard User Entity}

If the domain has no User-like entity, inject one:

\begin{lstlisting}[language=C]
fn standard_user_entity() -> EntityDef {
    EntityDef {
        name: "User".into(),
        fields: vec![
            FieldDef { name: "id".into(), rust_type: "i64".into(), sql_type: "BIGSERIAL PRIMARY KEY".into(), source: SystemGenerated, .. },
            FieldDef { name: "email".into(), rust_type: "String".into(), sql_type: "VARCHAR(255) UNIQUE NOT NULL".into(), source: UserInput, .. },
            FieldDef { name: "password_hash".into(), rust_type: "String".into(), sql_type: "VARCHAR(255) NOT NULL".into(), source: SystemGenerated, .. },
            FieldDef { name: "name".into(), rust_type: "String".into(), sql_type: "VARCHAR(100) NOT NULL".into(), source: UserInput, .. },
            FieldDef { name: "role".into(), rust_type: "String".into(), sql_type: "VARCHAR(30) NOT NULL DEFAULT 'viewer'".into(), source: UserInput, .. },
            FieldDef { name: "is_active".into(), rust_type: "bool".into(), sql_type: "BOOLEAN NOT NULL DEFAULT TRUE".into(), source: SystemGenerated, .. },
            FieldDef { name: "created_at".into(), rust_type: "String".into(), sql_type: "TIMESTAMPTZ NOT NULL DEFAULT NOW()".into(), source: SystemGenerated, .. },
            FieldDef { name: "updated_at".into(), rust_type: "String".into(), sql_type: "TIMESTAMPTZ NOT NULL DEFAULT NOW()".into(), source: SystemGenerated, .. },
        ],
        business_rules: vec![],
    }
}
\end{lstlisting}

\subsection{Prompt Enhancement: Entity Context from Norms}

When generating a service file, the prompt includes the entity's
business rules from the norm composition:

\begin{lstlisting}[language=C]
fn entity_context_from_norms(
    entity_name: &str,
    plan: &ComposedPlan,
) -> String {
    let mut ctx = String::new();

    // Business rules
    for rule in &plan.business_rules {
        if rule.affects_entity(entity_name) {
            ctx.push_str(&format!("- Business Rule: {}\n", rule.description));
        }
    }

    // Cross-entity wirings
    for wiring in &plan.wirings {
        if wiring.involves_entity(entity_name) {
            ctx.push_str(&format!("- Cross-Entity: {}\n", wiring.description));
        }
    }

    // State machine transitions (if applicable)
    if let Some(sm) = plan.state_machines.get(entity_name) {
        ctx.push_str(&format!("- State Machine states: {:?}\n", sm.states));
        ctx.push_str(&format!("- Valid transitions: {:?}\n", sm.transitions));
    }

    ctx
}
\end{lstlisting}

% ============================================================
\section{Multi-Domain Verification}

With \texttt{ForgePlan::from\_composed\_plan()}, any domain that
the norm system can compose should work. Verify:

\begin{lstlisting}[language=bash,title=Three Domains Must Work]
# Warehouse (via TOML)
isls forge-v2 --requirements examples/warehouse.toml \
    --mock-oracle --output ./output/warehouse
cd output/warehouse/backend && cargo build && cd ../../..

# E-Commerce (via TOML)
isls forge-v2 --requirements examples/ecommerce.toml \
    --mock-oracle --output ./output/ecommerce
cd output/ecommerce/backend && cargo build && cd ../../..

# Project Tracker (via TOML)
isls forge-v2 --requirements examples/project.toml \
    --mock-oracle --output ./output/project
cd output/project/backend && cargo build && cd ../../..

# Warehouse via Chat (with API key)
isls serve &
curl -X POST http://localhost:8420/api/chat \
  -H "Content-Type: application/json" \
  -d '{"message":"warehouse management system"}'
# Should return norms + plan_id + action buttons
\end{lstlisting}

% ============================================================
\section{CLI: isls serve with API Key}

\begin{lstlisting}[language=bash]
# Start Studio with mock mode (no LLM, thin generation)
isls serve

# Start Studio with API key (LLM generation available)
isls serve --api-key $OPENAI_API_KEY

# Or via environment
export OPENAI_API_KEY=sk-...
isls serve

# Custom port
isls serve --port 9090
\end{lstlisting}

The API key is passed to the AppState and used when
\texttt{/api/projects/forge} is called without a client-side key.

% ============================================================
\section{Implementation Plan}

\begin{enumerate}
\item \textbf{AppState with NormRegistry + plan storage.}
  Shared state in gateway with Mutex-wrapped registries.
  Test: start serve, GET /api/norms returns 24 norms.

\item \textbf{POST /api/chat --- real intent extraction + norm matching.}
  Connect isls-chat + isls-norms to the handler.
  Test: POST ``warehouse management'' returns activated norms.

\item \textbf{POST /api/projects/forge --- real generation.}
  Connect isls-forge-llm to the handler.
  Test: forge with mock-oracle produces files on disk.

\item \textbf{ForgePlan::from\_composed\_plan() --- dynamic plans.}
  Replace warehouse\_default() with norm-driven plan building.
  Test: ecommerce norms produce ecommerce ForgePlan.

\item \textbf{Standard User entity injection.}
  If domain has no User, inject one.
  Test: ecommerce (has Customer, no User) gets User injected.

\item \textbf{Studio chat wiring.}
  JavaScript: sendMessage(), handleAction(), loadProjects().
  Test: type in browser, see response, click Generate.

\item \textbf{Studio file browser.}
  GET /api/projects/:id/files, file viewer panel.
  Test: click project, browse files, view content.

\item \textbf{Studio norm explorer wiring.}
  Highlight activated norms when chat response includes them.
  Test: norms light up when user describes an app.

\item \textbf{Multi-domain verification.}
  All three TOML domains compile in mock mode.
  Test: warehouse + ecommerce + project = 0 errors each.

\item \textbf{End-to-end: chat in browser $\to$ working app.}
  Type ``warehouse management'' $\to$ Generate $\to$ files appear $\to$ compilable.
  Test: the complete flow in one browser session.
\end{enumerate}

% ============================================================
\section{Success Criteria}

\begin{enumerate}
\item \texttt{cargo build --workspace} --- zero errors.
\item All existing 340+ tests pass.
\item \texttt{isls serve} starts and Studio loads at :8420/studio.
\item POST /api/chat with ``warehouse'' returns norms + plan.
\item POST /api/projects/forge generates files on disk.
\item Studio chat: type description $\to$ see norms $\to$ click Generate
  $\to$ project appears.
\item Studio file browser: browse and view generated files.
\item Norm explorer highlights activated norms per chat message.
\item \texttt{forge-v2 --mock-oracle} produces 0-error backends for
  warehouse, ecommerce, and project tracker.
\item With \texttt{--api-key}: warehouse backend compiles after
  LLM generation with cargo-check retry loop.
\item ForgePlan is built dynamically from ComposedPlan, not
  hardcoded per domain.
\item E-Commerce gets a User entity injected automatically.
\item All generated code is template-free (no Tera artifacts).
\end{enumerate}

% ============================================================
\section{Constraints for Claude Code}

\begin{enumerate}
\item \textbf{Repository: graphity/isls.}
\item \textbf{No new crates.} All work in existing crates.
\item \textbf{v3.1 functionality must not break.} CLI forge-v2
  still works independently of the gateway.
\item \textbf{The gateway is synchronous for v3.2.} No async
  streaming. The forge call blocks and returns when done.
  Progress is polled via GET /api/projects/:id if needed.
\item \textbf{Studio is still ONE HTML file.} All JavaScript
  inline. No npm. No build step.
\item \textbf{API key handling:} accept via --api-key flag,
  OPENAI\_API\_KEY env var, or per-request in forge body.
  Never log or persist the key.
\item \textbf{AppState uses Arc$<$Mutex$<$T$>>$} for shared data.
  The gateway runs single-threaded for v3.2.
\item \textbf{Sequential:} AppState $\to$ chat handler $\to$
  forge handler $\to$ dynamic ForgePlan $\to$ user injection $\to$
  studio JS $\to$ file browser $\to$ norm explorer $\to$
  multi-domain verify $\to$ end-to-end test.
\item \textbf{The chat-to-generation flow is the most important
  deliverable.} If time runs out: skip file browser, skip
  norm explorer animation. But chat $\to$ generate must work.
\item \textbf{No \texttt{unwrap()} in library code.}
\item \textbf{Every public item documented.}
\item \textbf{MIT License. Copyright (c) 2026 Sebastian Klemm.}
\end{enumerate}

\vfill
\noindent\rule{\textwidth}{0.4pt}
\begin{center}
\small
ISLS v3.2 --- Full Pipeline Wiring --- Sebastian Klemm --- March 2026\\
Type in the browser. Get a working app.\\
All pieces exist. Now they talk to each other.\\
\url{https://github.com/LashSesh/graphity}
\end{center}

\end{document}
